Write $u = \tfrac12 - \delta$, $\lambda = 2\pi/r$, $s(x) = 1/(e^{x} + 1)$, and let $\mathcal{A}$ be the set of the numbers $k + \delta$ and $k + 1 - \delta$, $k \ge 0$, so $T = -r^{2}/24 + \sum_{a \in \mathcal{A}}a\,s(\lambda a)$. On the conic $B_{2}(\delta) = (\delta - \tfrac12)^{2} - \tfrac{1}{12} = -r^{2}/12$, so the sea energy $\tfrac{\kappa}{2r}B_{2}(\delta)$ of \claimref{synthesis.null-ordering-of-counted-units-is-zeta-symmetric}(ii) is $-\kappa r/24$, and the mean energy in the Gibbs state at $\beta = 2\pi/\kappa$ is that plus $\sum_{a}(\kappa/r)a\,s(2\pi a/r)$: it is $(\kappa/r)T$. The function is symmetric under $\delta \mapsto 1 - \delta$, which exchanges the two families of $\mathcal{A}$ and fixes $r$; it is real analytic on the open interval of the statement, the series converging uniformly with its derivatives on compact subsets (Step 2).

\emph{Step 1: no zero with $r \le \tfrac{3}{20}$.} For $\delta$ in the interval, $\delta > \tfrac12 - \tfrac{1}{\sqrt{12}} > \tfrac15$ (because $(\tfrac12 - \tfrac15)^{2} = \tfrac{9}{100} > \tfrac{1}{12}$) and $1 - \delta \ge \tfrac12$, so the two elements of $\mathcal{A}$ with index $k$ are at least $k + \tfrac15$. For $r \le \tfrac{3}{20}$, $\lambda \ge 40\pi/3 > 8$; the function $a \mapsto ae^{-\lambda a}$ decreases for $a \ge 1/\lambda$, in particular for $a \ge \tfrac15$. With $s(x) \le e^{-x}$,
\[
\sum_{a \in \mathcal{A}}a\,s(\lambda a) \le 2\sum_{k \ge 0}(k + \tfrac15)e^{-\lambda(k + 1/5)} \le 2e^{-\lambda/5}\Bigl[\tfrac15 + \tfrac65\,\frac{e^{-\lambda}}{(1 - e^{-\lambda})^{2}}\Bigr] < \tfrac{41}{100}\,e^{-2\pi/(5r)},
\]
using $(k + \tfrac15) \le \tfrac65 k$ for $k \ge 1$, $\sum_{k \ge 1}ke^{-\lambda k} = e^{-\lambda}/(1 - e^{-\lambda})^{2}$ and $e^{-8} < \tfrac{1}{2000}$. So $T < 0$ as soon as $r^{2}e^{2\pi/(5r)} > 24 \cdot \tfrac{41}{100}$. The function $r^{2}e^{c/r}$, $c = 2\pi/5$, has derivative $e^{c/r}(2r - c)$, negative for $r < \pi/5$; so on $0 < r \le \tfrac{3}{20} < \tfrac{\pi}{5}$ it is at least its value at $\tfrac{3}{20}$, which is $\tfrac{9}{400}e^{8\pi/3} \ge \tfrac{9}{400}e^{8} > \tfrac{9}{400}(2.718)^{8} > 66 > 24 \cdot \tfrac{41}{100}$ (the check verifies $e > 2.718$, $\pi \ge 3$ and the last inequality in rationals). On the conic $r$ increases with $\delta$ for $\delta \le \tfrac12$, and at $\delta = \tfrac{429}{2000}$, $r^{2} = 1 - 12\bigl(\tfrac{571}{2000}\bigr)^{2} < \bigl(\tfrac{3}{20}\bigr)^{2}$. So $T < 0$ for $\tfrac12 - \tfrac{1}{\sqrt{12}} < \delta \le \tfrac{429}{2000}$.

\emph{Step 2: the omitted terms.} The check sums the terms with $k < 8$ and bounds the rest. A term is $g = a\,s(x)$ with $x = \lambda a$, $a' = \pm 1$, $a'' = 0$ (derivatives in $\delta$). For $x \ge 0$: $0 < s \le e^{-x}$, $|s_{x}| = s(1 - s) \le e^{-x}$, $|s_{xx}| = s(1 - s)|1 - 2s| \le e^{-x}$. Let $M \ge \max(1, \lambda, |\lambda'|, |\lambda''|)$ on the interval of $\delta$ considered. Then $|x'| = |\lambda'a + \lambda a'| \le M(a + 1)$, $|x''| = |\lambda''a + 2\lambda'a'| \le M(a + 2)$, and
\[
|g| \le ae^{-x}, \qquad |g'| \le e^{-x}\bigl(1 + aM(a + 1)\bigr), \qquad |g''| \le e^{-x}\bigl(2M(a + 1) + aM^{2}(a + 1)^{2} + aM(a + 2)\bigr).
\]
For $a \ge 2$ each right side is at most $4M^{2}a^{3}e^{-x}$, hence at most $4M^{2}a^{3}e^{-\lambda a}$; the weaker $11M^{2}a^{3}e^{-\lambda a}$ is used, which gives the constant below. The terms omitted are, for each $k \ge 8$, two values of $a$ in $[k, k + 1]$, so each of the three sums is at most $22M^{2}\sum_{k \ge 8}(k + 1)^{3}\rho^{k}$ with $\rho = e^{-\lambda_{0}}$, $\lambda_{0}$ a lower bound of $\lambda$ on the interval. Since $\lambda \ge 2\pi$, $\rho < \tfrac{1}{500}$, the ratio of consecutive terms is at most $(\tfrac{10}{9})^{3}\rho < \tfrac{1}{300}$, and the sum is at most $\tfrac{300}{299}\,9^{3}\rho^{8}$. So each of the three omitted sums is bounded by $\tfrac{223}{10}\,M^{2}\,9^{3}\,e^{-8\lambda_{0}}$, the amount by which the check widens its enclosures of $T$, $T'$ and $T''$, with $M$ and $\lambda_{0}$ read from its enclosure of $\lambda$ and its derivatives. The same bound, applied on compact sets, gives the uniform convergence used above.

\emph{Step 3: the enclosures.} An interval is a pair of integers standing for $[\mathrm{lo}/2^{128}, \mathrm{hi}/2^{128}]$; sums, products, reciprocals and square roots round the lower end down and the upper end up; the exponential at an end $t$ is computed by halving $t$ until $|t| \le \tfrac12$, summing its series to order $30$ with interval terms (the remainder is below $2^{-140}$ and is covered by widening the result by one unit), and squaring back; $\pi = 16\arctan\tfrac15 - 4\arctan\tfrac{1}{239}$ is enclosed by consecutive partial sums of the two alternating series. A function of $\delta$ on an interval is carried as three intervals enclosing it and its first two derivatives, combined by the rules for sums, products, reciprocals, square roots and for $s$ (with $s' = -s(1 - s)$, $s'' = s(1 - s)(1 - 2s)$ and the chain rule). Each rule is valid for enclosures because the operations are monotone in the inclusion of intervals. So for a rational interval $[\delta_{1}, \delta_{2}]$ inside the domain the check returns intervals containing every value of $T$, $T'$ and $T''$ on it, after the widening of Step 2.

\emph{Step 4: signs on three stretches.} For an interval with midpoint $\delta_{m}$ and half-width $h$, the mean value theorem gives $T(\delta) \in T(\delta_{m}) + T'([\delta_{1}, \delta_{2}])\cdot[-h, h]$. The check evaluates the right side, bisects when its sign is undecided, and finds (73 intervals in all): $T < 0$ on $[\tfrac{429}{2000}, \tfrac{246}{1000}]$, $T > 0$ on $[\tfrac{254}{1000}, \tfrac{3233}{10000}]$, $T < 0$ on $[\tfrac{3433}{10000}, \tfrac{46}{100}]$.

\emph{Step 5: near the three points.} The check finds $T' > 0$ on $[\tfrac{246}{1000}, \tfrac{254}{1000}]$ (16 pieces), $T' < 0$ on $[\tfrac{3233}{10000}, \tfrac{3433}{10000}]$ (20 pieces) and $T'' < 0$ on $[\tfrac{46}{100}, \tfrac12]$ (40 pieces).

\emph{Step 6: the zeros.} By clause (ii) of \claimref{synthesis.window-energy-tests-with-free-spacing-admit-twists-half-third-and-quarter} the thermal test holds on the conic at twists $\tfrac14$, $\tfrac13$, $\tfrac12$, so $T$ vanishes there. On $[\tfrac{246}{1000}, \tfrac{254}{1000}]$, $T$ is strictly increasing, so $\tfrac14$ is its only zero there and it is simple; likewise $\tfrac13$ on $[\tfrac{3233}{10000}, \tfrac{3433}{10000}]$, where $T$ is strictly decreasing. By the symmetry $T(\delta) = T(1 - \delta)$, $T'(\tfrac12) = 0$. On $[\tfrac{46}{100}, \tfrac12]$, $T'' < 0$, so $T'$ is strictly decreasing and $T'(\delta) > T'(\tfrac12) = 0$ for $\delta < \tfrac12$; hence $T$ is strictly increasing there and $T(\delta) < T(\tfrac12) = 0$ for $\tfrac{46}{100} \le \delta < \tfrac12$, and the zero at $\tfrac12$ is double ($T''(\tfrac12) < 0$). With Steps 1 and 4 the intervals cover $(\tfrac12 - \tfrac{1}{\sqrt{12}}, \tfrac12]$ and give the signs stated. The conic has no point with $r > 1$, every point of it with $r \le 1$ has twist $\delta$ or $1 - \delta$ with $\delta$ in the interval treated, and the ground test holds exactly on the conic (clause (i) of the cited theorem); so the points passing both tests are the five listed.
